\section*{Bilan du projet}

Ce projet de fin d'études nous a permis de concevoir et de développer
\textbf{TAC-TIC IntelliPM}, une plateforme de gestion de projets complète,
intelligente et enterprise-ready. En \textbf{six sprints Scrum} de deux semaines
chacun, nous avons livré l'intégralité des \textbf{34 user stories} planifiées,
représentant \textbf{228 story points}.

\begin{table}[H]
  \centering
  \caption{Récapitulatif des six sprints et leur vélocité}
  \label{tab:sprints_recap}
  \begin{tabularx}{\textwidth}{C{1.5cm}L{6cm}C{2cm}C{2cm}}
    \toprule
    \thead{Sprint} & \thead{Thème principal} & \thead{US livrées} & \thead{Points} \\
    \midrule
    \rowcolor{lightgray}
    1 & Gestion des utilisateurs et authentification & 5  & 26 \\
    2 & Projets, équipes et invitations              & 6  & 35 \\
    \rowcolor{lightgray}
    3 & Tâches, Kanban et dépendances                & 5  & 39 \\
    4 & Planification visuelle et notifications       & 8  & 52 \\
    \rowcolor{lightgray}
    5 & Intelligence artificielle hybride            & 4  & 34 \\
    6 & Sécurité enterprise et journal d'audit        & 6  & 42 \\
    \midrule
    \multicolumn{2}{r}{\textbf{Total}} & \textbf{34} & \textbf{228} \\
    \bottomrule
  \end{tabularx}
\end{table}

\medskip
Les principaux résultats techniques obtenus sont les suivants :

\begin{itemize}
  \item \textbf{Architecture distribuée opérationnelle} : backend Laravel~12
        exposant plus de 70 endpoints REST groupés en 38 routes, SPA React~19
        intégralement en français, et microservice FastAPI Python~3.11 pour les
        modèles ML.
  \item \textbf{Sécurité enterprise multicouche} : authentification Sanctum
        (Bearer Token), RBAC à trois niveaux (système + projet + permissions
        atomiques), Policies Laravel, journal d'audit enrichi
        (valeurs avant/après, IP, sévérité automatique, export PDF). Protection
        OWASP Top~10 couverte.
  \item \textbf{Temps réel complet} : 11 événements WebSocket broadcastés via
        Laravel Reverb, tous couverts par des listeners frontend dédiés — latence
        médiane inférieure à 20\,ms.
  \item \textbf{IA opérationnelle} : AUC ROC 0,92 (détection de risques),
        MAE 1,42\,h (estimation de durée), inférence ML $<$ 45\,ms (P50).
        Assistant conversationnel LLM (Llama~3.3 70B via OpenRouter) avec
        contexte projet chargé en temps réel.
  \item \textbf{Analytics avancées} : tableau de bord d'audit avec détection
        d'anomalies (tâches instables, membres inactifs, pics d'activité),
        graphiques CSS/SVG pur — zéro dépendance externe.
  \item \textbf{Qualité certifiée} : SonarCloud Quality Gate passé ---
        0 bug critique, 0 vulnérabilité, duplication $\approx$ 2\,\%.
  \item \textbf{Performance} : P95 API $<$ 500\,ms, temps de chargement initial
        SPA $<$ 2\,s (code splitting + TanStack Query). Élimination des requêtes
        N+1 critiques (Gantt, AiController).
\end{itemize}

\section*{Déploiement}

L'application est conçue pour un déploiement cloud en trois composants indépendants :

\begin{table}[H]
  \centering
  \begin{tabularx}{\textwidth}{L{3cm}L{4cm}X}
    \toprule
    \textbf{Composant} & \textbf{Infrastructure} & \textbf{Détails} \\
    \midrule
    \rowcolor{lightgray}
    Frontend React  & Vercel (CDN mondial)    & Build Vite, déploiement automatique GitHub Actions \\
    Backend Laravel & VPS 4 vCPU / 8\,Go RAM & PHP-FPM + Nginx, Redis, Laravel Horizon (queues) \\
    \rowcolor{lightgray}
    ML FastAPI      & VPS 2 vCPU / 4\,Go RAM & Uvicorn, modèles persistés Joblib, rechargement à chaud \\
    Base de données & OVH MySQL 8 managé      & Sauvegardes automatiques, réplication lecture \\
    \bottomrule
  \end{tabularx}
\end{table}

La chaîne CI/CD (GitHub Actions) orchestre : lint, tests PHPUnit, analyse SonarCloud,
build Docker, déploiement automatique sur branche \texttt{main}. Cloudflare assure le
TLS, le CDN et la protection DDoS.

\section*{Apports personnels}

Ce projet nous a offert une formation transversale couvrant plusieurs domaines :

\begin{description}[leftmargin=3cm]
  \item[Génie logiciel] Architecture en couches, patterns Repositories et Services,
        SOLID, Clean Architecture.
  \item[Bases de données] Modélisation ERD, optimisation des requêtes Eloquent,
        migrations, indexation.
  \item[Intelligence artificielle] Préparation de données, feature engineering,
        sélection de modèles, validation croisée.
  \item[Sécurité] RBAC, politiques d'autorisation, protection OWASP, audit trail.
  \item[Méthodes agiles] Cérémonies Scrum, story points, vélocité, rétrospective.
  \item[DevOps] CI/CD GitHub Actions, Docker Compose, déploiement cloud.
\end{description}

\section*{Limites}

Nous identifions les principales limites de la version actuelle :

\begin{itemize}
  \item Les données d'entraînement des modèles ML sont synthétiques (simulées) ;
        en production réelle, l'historique de projets alimenterait des modèles
        plus précis et permettrait un réentraînement continu.
  \item L'application est uniquement disponible en version web ; une application
        mobile (React Native) permettrait un accès en mobilité.
  \item La fonctionnalité d'intégration GitHub/GitLab (lier des commits à des
        tâches) a été reportée faute de temps.
  \item Le système d'export PDF, bien que fonctionnel, repose sur
        \texttt{window.print()} ; une génération serveur via Puppeteer ou
        wkhtmltopdf offrirait des rapports plus robustes.
\end{itemize}

\section*{Perspectives}

Les évolutions envisagées pour TAC-TIC IntelliPM s'articulent autour de trois axes :

\paragraph{Enrichissement fonctionnel}
Application mobile React Native, tableau de bord exécutif consolidé, chat intégré
aux tâches, localisation arabe (RTL), gamification (badges, points de productivité),
export Excel du Gantt, génération PDF serveur (Puppeteer), intégration GitHub/GitLab.

\paragraph{Intelligence artificielle avancée}
Réentraînement continu sur les données réelles (pipeline MLOps), analyse NLP des
commentaires pour détection de conflits, valeurs SHAP pour l'explicabilité des
recommandations, détection automatique de surcharge d'équipe, modèle de prédiction
de vélocité de sprint.

\paragraph{Industrialisation SaaS}
Déploiement Kubernetes/Helm, architecture multi-tenant avec isolation des données
par organisation, connecteurs Teams/Slack, certification ISO~27001, audit trail
immuable (blockchain légère), modèle de facturation par siège.

\bigskip
\noindent\rule{\linewidth}{0.5pt}

\medskip
Ce projet a été l'occasion de cultiver l'équilibre entre rigueur technique
et humilité intellectuelle, entre précision de l'ingénieur et curiosité du
chercheur. Nous espérons que TAC-TIC IntelliPM constituera une base solide
pour une suite industrielle au service des équipes de développement tunisiennes
et au-delà.
